\documentclass[11pt]{article}

\usepackage{amsmath,amssymb,amsthm,mathtools}
\usepackage[margin=1in]{geometry}
\usepackage{microtype}

\newtheorem{theorem}{Theorem}[section]
\newtheorem{lemma}[theorem]{Lemma}
\newtheorem{corollary}[theorem]{Corollary}
\newtheorem{remark}[theorem]{Remark}

\newcommand{\dd}{\,\mathrm d}

\title{Pendant Leaves on the Outer Orbit of a Bowtie:\\
Double Symmetrization and Fractional Triangle Bias}
\author{}
\date{}

\begin{document}

\maketitle

\begin{abstract}
Let $B$ be the bowtie graph, the vertex amalgam of two triangles.  Its four
noncentral vertices form one automorphism orbit.  We attach $h$ pendant
leaves arbitrarily among those four vertices and prove the
chromatic-polynomial graphon bound for $1\leq h\leq6$ and every
$p\geq1/2$.  Averaging over the outer orbit gives degree exponent $h/4$ at
each outer vertex.  Conditioning on the shared center then factorizes the
density as the integral of a square.  Cauchy--Schwarz and a second
symmetrization inside a triangle produce the uniform fractional weight
$d^{h/6}$, which is controlled by a biased-measure H\"older argument and the
accepted triangle bound.  The case $h=1$ proves the formerly open NetworkX
Atlas graph $119$.
\end{abstract}


%%%%%%%%%%%%%%%%%%%%%%%%%%%%%%%%%%%%%%%%%%%%%%%%%%%%%%%%%%%%%%%%%%%%%%%%%%%%%%%
\section{Bowtie family and exact target}
%%%%%%%%%%%%%%%%%%%%%%%%%%%%%%%%%%%%%%%%%%%%%%%%%%%%%%%%%%%%%%%%%%%%%%%%%%%%%%%

Let $(\Omega,\mu)$ be an arbitrary probability space and let
$W\colon\Omega^2\to[0,1]$ be symmetric and measurable.  Write
\[
 p=\int_{\Omega^2}W(x,y)\dd\mu(x)\dd\mu(y),
 \qquad
 d(x)=\int_\Omega W(x,y)\dd\mu(y).
\]

Let $B$ have center $o$ and outer vertices $a,b,c,d$, with triangles
$oab$ and $ocd$.  The automorphism group fixes $o$ and is transitive on
$\{a,b,c,d\}$: it may swap the two triangles and independently swap the
endpoints within either outer edge.

For nonnegative integers $n_a,n_b,n_c,n_d$, put
\[
 h=n_a+n_b+n_c+n_d
\]
and let $B(n_a,n_b,n_c,n_d)$ be obtained by attaching $n_v$ private
pendant leaves to outer vertex $v$.  We assume
\[
 1\leq h\leq6.
\]

The bowtie is the vertex amalgam of two triangles, so
\[
 \chi_B(x)=\frac{\chi_{K_3}(x)^2}{x}
 =x(x-1)^2(x-2)^2.
\]
Each leaf supplies a further factor $x-1$.  Therefore
\begin{equation}
 \chi_{B(\boldsymbol n)}(x)
 =x(x-1)^{h+2}(x-2)^2,
 \qquad \chi(B(\boldsymbol n))=3.
 \label{eq:chromatic}
\end{equation}
If
\[
 A_3(p):=p(2p-1),
\]
then direct substitution in \eqref{eq:chromatic} gives
\begin{equation}
 \Phi_{B(\boldsymbol n)}(p)
 =p^hA_3(p)^2
 =p^{h+2}(2p-1)^2.
 \label{eq:target}
\end{equation}
This includes the polynomial continuation at $p=1$; at the required
threshold $p=1/2$, it vanishes.

We use the accepted triangle bound in the form
\begin{equation}
 t(K_3,V)\geq A_3(s)=s(2s-1)
 \label{eq:triangle-input}
\end{equation}
for every graphon $V$ of edge density $s\geq1/2$, on an arbitrary
probability space.


%%%%%%%%%%%%%%%%%%%%%%%%%%%%%%%%%%%%%%%%%%%%%%%%%%%%%%%%%%%%%%%%%%%%%%%%%%%%%%%
\section{Fractional triangle bias}
%%%%%%%%%%%%%%%%%%%%%%%%%%%%%%%%%%%%%%%%%%%%%%%%%%%%%%%%%%%%%%%%%%%%%%%%%%%%%%%

\begin{lemma}[Fractionally degree-weighted edge]
\label{lem:fractional-edge}
For $0<\beta\leq1$, define
\[
 N_\beta=\int_{\Omega^2}W(x,y)d(x)^\beta d(y)^\beta
 \dd\mu(x)\dd\mu(y).
\]
Then
\[
 N_\beta\geq p^{1+2\beta}.
\]
\end{lemma}

\begin{proof}
On $\{d(x)=0\}$, define $W(x,y)/d(x)$ to be zero.  Fubini implies that
$W(x,y)=0$ for almost every $y$ there.  Put
\[
 \theta=\frac1{1+2\beta},
 \qquad
 \eta=\frac\beta{1+2\beta}.
\]
Then $\theta+2\eta=1$, and weighted H\"older applied to
\[
 W
 =\bigl(Wd(x)^\beta d(y)^\beta\bigr)^\theta
  \left(\frac W{d(x)}\right)^\eta
  \left(\frac W{d(y)}\right)^\eta
\]
gives
\[
 p\leq N_\beta^\theta
 \left(\int\frac W{d(x)}\dd\mu^2\right)^\eta
 \left(\int\frac W{d(y)}\dd\mu^2\right)^\eta.
\]
The quotient integrals are at most one, proving the claim.
\end{proof}

\begin{lemma}[Triangle rescaling]
\label{lem:scalar}
Let $0<h\leq6$, let $p\in(1/2,1]$, and suppose
$0<z\leq p^{h/2}$.  Put
\[
 s_0=p\left(\frac{p^{h/2}}z\right)^{2/3}.
\]
If $s_0\leq1$, then
\[
 zA_3(s_0)\geq p^{h/2}A_3(p).
\]
\end{lemma}

\begin{proof}
For $s\in[1/2,1]$, set
\[
 \psi(s)=\frac{A_3(s)}{s^{3/2}}
 =2\sqrt{s}-\frac1{\sqrt{s}}.
\]
Then
\[
 \psi'(s)=\frac1{\sqrt{s}}+\frac1{2s^{3/2}}>0.
\]
The definition of $s_0$ is equivalent to
\[
 z=p^{h/2}\left(\frac p{s_0}\right)^{3/2}.
\]
Since $s_0\geq p$, monotonicity gives
\[
 zA_3(s_0)
 =p^{(h+3)/2}\psi(s_0)
 \geq p^{(h+3)/2}\psi(p)
 =p^{h/2}A_3(p).
\]
\end{proof}

The function $A_3(s)=s(2s-1)$ is nondecreasing on $[1/2,1]$.


%%%%%%%%%%%%%%%%%%%%%%%%%%%%%%%%%%%%%%%%%%%%%%%%%%%%%%%%%%%%%%%%%%%%%%%%%%%%%%%
\section{The outer-leaf theorem}
%%%%%%%%%%%%%%%%%%%%%%%%%%%%%%%%%%%%%%%%%%%%%%%%%%%%%%%%%%%%%%%%%%%%%%%%%%%%%%%

\begin{theorem}[Leaves on the bowtie's outer orbit]
\label{thm:main}
Let $1\leq h\leq6$ and let
$n_a+n_b+n_c+n_d=h$.  Every graphon $W$ of edge density
$p\in[1/2,1]$ satisfies
\[
 \boxed{
 t\bigl(B(n_a,n_b,n_c,n_d),W\bigr)
 \geq p^hA_3(p)^2
 =\Phi_{B(\boldsymbol n)}(p).
 }
\]
\end{theorem}

\begin{proof}
The threshold $p=1/2$ follows from nonnegativity, and at $p=1$ both sides
equal one.  Assume $p\in(1/2,1)$.

Integrating the leaves first gives the bowtie core integrand multiplied by
\[
 d(x_a)^{n_a}d(x_b)^{n_b}d(x_c)^{n_c}d(x_d)^{n_d}.
\]
Average over the automorphism group of the bowtie.  Its action is transitive
on the four outer vertices, so each exponent has group average $h/4$.
The arithmetic--geometric mean inequality therefore gives the lower bound
in which every outer variable has exponent
\[
 \alpha=\frac h4.
\]

For a fixed image $x$ of the center, define
\[
 R_\alpha(x)
 =\int_{\Omega^2}
 W(x,y)W(x,z)W(y,z)d(y)^\alpha d(z)^\alpha
 \dd\mu(y)\dd\mu(z).
\]
The two outer edges $ab$ and $cd$ use disjoint variables.  Hence the
symmetrized lower bound factorizes exactly as
\begin{equation}
 t(B(\boldsymbol n),W)
 \geq\int_\Omega R_\alpha(x)^2\dd\mu(x)
 \geq\left(\int_\Omega R_\alpha(x)\dd\mu(x)\right)^2,
 \label{eq:square}
\end{equation}
where the second inequality is Cauchy--Schwarz.

The last integral is a triangle density with weight $d^\alpha$ at two of
its three vertices.  By triangle symmetry, the three choices of unweighted
vertex give equal integrals.  Average those three representations and apply
the arithmetic--geometric mean inequality again.  Every triangle vertex then
has exponent
\[
 \beta=\frac{2\alpha}{3}=\frac h6\in(0,1].
\]
Consequently
\begin{equation}
 \int R_\alpha
 \geq\int_{\Omega^3}
 W(x,y)W(y,z)W(z,x)
 d(x)^\beta d(y)^\beta d(z)^\beta\dd\mu^3.
 \label{eq:triangle-symmetry}
\end{equation}

Put
\[
 M=\int_\Omega d(x)^\beta\dd\mu(x),
 \qquad z=M^3,
 \qquad
 \dd\nu(x)=\frac{d(x)^\beta}{M}\dd\mu(x).
\]
Here $M>0$ because $p>0$.  The right side of
\eqref{eq:triangle-symmetry} is $z\,t(K_3,W;\nu)$.  The edge density on
$(\Omega,\nu)$ is
\[
 s=\frac{N_\beta}{M^2}.
\]
Concavity of $u^\beta$ and Lemma~\ref{lem:fractional-edge} give
\[
 0<z\leq p^{3\beta}=p^{h/2},
\]
and
\[
 s\geq\frac{p^{1+2\beta}}{M^2}
 =p\left(\frac{p^{h/2}}z\right)^{2/3}
 =:s_0.
\]
Since $s\leq1$, also $s_0\leq1$; moreover $s\geq s_0\geq p>1/2$.

Apply the accepted triangle bound \eqref{eq:triangle-input}, its
monotonicity, and Lemma~\ref{lem:scalar}:
\[
 \int R_\alpha
 \geq zA_3(s)
 \geq zA_3(s_0)
 \geq p^{h/2}A_3(p).
\]
Substitution in \eqref{eq:square} gives
\[
 t(B(\boldsymbol n),W)
 \geq p^hA_3(p)^2,
\]
which is \eqref{eq:target}.
\end{proof}

Every balanced complete $\ell$-partite graphon, $\ell\geq2$, realizes
equality.  Its degree is constant, the rooted weighted-triangle function is
constant by class symmetry, and the accepted triangle bound is sharp.  At
$\ell=2$, both sides vanish.


%%%%%%%%%%%%%%%%%%%%%%%%%%%%%%%%%%%%%%%%%%%%%%%%%%%%%%%%%%%%%%%%%%%%%%%%%%%%%%%
\section{Atlas 119, formalization audit, and limitations}
%%%%%%%%%%%%%%%%%%%%%%%%%%%%%%%%%%%%%%%%%%%%%%%%%%%%%%%%%%%%%%%%%%%%%%%%%%%%%%%

\begin{corollary}[Atlas 119]\label{cor:atlas119}
NetworkX Atlas graph $119$, with graph6 code \texttt{EtoO}, is positive.  It
has order six, size seven, chromatic number three, and
\[
 \chi_{F_{119}}(x)=x(x-1)^3(x-2)^2.
\]
Every graphon of edge density $p\in[1/2,1]$ satisfies
\[
 t(F_{119},W)\geq p^3(2p-1)^2=\Phi_{F_{119}}(p).
\]
\end{corollary}

\begin{proof}
Its NetworkX edge list is
\[
 01,\ 02,\ 03,\ 04,\ 14,\ 23,\ 35.
\]
The triangles $014$ and $023$ share center $0$, and vertex $5$ is a leaf at
the outer vertex $3$.  Thus Theorem~\ref{thm:main} applies with $h=1$.
\end{proof}

All uses of Tonelli--Fubini concern bounded nonnegative measurable
integrands.  The proof explicitly handles $\{d=0\}$ and proves the biased
normalization positive before dividing.  It uses ordinary homomorphism
densities and assumes no regularity, finite-step model, minimum degree, or
injectivity.  The only accepted density input is the triangle theorem with
exactly matching edge-density range.

The leaves must attach only to the four-vertex outer orbit.  The shared
center carries no leaf in this theorem.  The bound $h\leq6$ is the exact
limit of this proof, because the second symmetrization produces exponent
$h/6$ and the needed concavity reverses beyond one.  The result does not
cover arbitrary leaves on a vertex-amalgam, tails of length two, or a general
mixed-chordal graph.

The graph identities, both symmetrized exponents, scalar derivative, and
Atlas identification are reconstructed exactly by
\begin{verbatim}
python codes/verify_bowtie_outer_leaves.py
\end{verbatim}
For formal verification, the new pieces are the outer-orbit
arithmetic--geometric symmetrization, conditional square factorization,
Cauchy--Schwarz step, and the second triangle symmetrization.  The
graph-specific module awaiting inspection is
\texttt{lean/Taeyoung/Examples/Graph119.lean}.

%%%%%%%%%%%%%%%%%%%%%%%%%%%%%%%%%%%%%%%%%%%%%%%%%%%%%%%%%%%%%%%%%%%%%%%%%%%%%%%
\section{What the Lean formalization actually proves}
%%%%%%%%%%%%%%%%%%%%%%%%%%%%%%%%%%%%%%%%%%%%%%%%%%%%%%%%%%%%%%%%%%%%%%%%%%%%%%%

The machine-checked development is \texttt{lean/Taeyoung/Methods/BowtieLeaf.lean},
ending at \texttt{satisfiesLowerBound\_119}.

\paragraph{Scope: $h=1$ only.}
Theorem~\ref{thm:main} is \emph{not} formalized for $1\leq h\leq6$.  What is
proved is Corollary~\ref{cor:atlas119}, on all of $p\in[1/2,1]$, and it is the
only member of the family inside the six-vertex catalogue: the bowtie already
has five vertices, so $h=1$.  The two symmetrisation exponents are therefore
$\alpha=1/4$ and $\beta=1/6$, both \emph{reciprocal integers}, which is what
lets the pre-existing \texttt{Methods/OddLeaf/Bias.lean} apply unchanged at
$m=6$: it builds $d^{1/m}$, $M=\int d^{1/m}$, the tilted measure $\nu$ and the
reread kernel for every $m\geq1$, and both $(d^{1/m})^m=d$ and $M^m\leq p$ stay
$\mathbb{N}$-powers.  A general $h$ would need $d^{h/6}$ with $h\nmid6$, hence
genuine \texttt{Real.rpow} arithmetic in those two places.

\paragraph{Lemma~\ref{lem:fractional-edge} was already available.}
It is the same statement as \texttt{odd\_cycle\_one\_leaf.tex} Lemma~2.1 at
$\beta=1/m$, and is formalized once, for every $m$, as
\texttt{OddLeaf.pow\_le\_pow\_rootEdge}:
\[
 p^{m+2}\leq N^m .
\]
The polynomial form is deliberate --- no real power of $p$ appears in any
statement.  Atlas 119 uses the $m=6$ instance
\texttt{OddLeaf.pow\_eight\_le\_pow\_rootEdge}, $p^8\leq N^6$.

\paragraph{The departure: the rescaling carries no derivative and no
square root.}
Lemma~\ref{lem:scalar} is proved here by differentiating
$\psi(s)=A_3(s)s^{-3/2}=2\sqrt s-1/\sqrt s$, and at $h=1$ its statement contains
$p^{h/2}=\sqrt p$.  The formalization forms no square root at all.  The
observation is that the quantity the proof actually feeds into \eqref{eq:square}
is the \emph{square} $(\int R_\alpha)^2$, so with $K=M^6=z^2$ the whole
rescaling can be stated squared,
\[
 p\,A_3(p)^2\leq K\,A_3(s)^2
 \qquad(\texttt{BowtieLeaf.transfer}),
\]
and its hypotheses are the two polynomial facts $K\leq p$ (Jensen,
\texttt{OddLeaf.pow\_rootMean\_le}) and $p^8\leq K^2s^6$ (the fractional
Hölder, through $N=M^2s$).  Clearing $s^3$ and dividing by $p^2s^2$ reduces it
to
\begin{equation}
 p(2s-1)^2-s(2p-1)^2=(s-p)(4ps-1)\geq0
 \qquad\bigl(\tfrac12\leq p\leq s\bigr),
 \label{eq:lean-bowtie}
\end{equation}
a single \texttt{ring} factorisation --- the same factor as the rescaling in
\texttt{Methods/CliqueDist/Diamond.lean}.  The intermediate $s_0$ never appears,
and neither does the hypothesis $s_0\leq1$; the inequality $p\leq s$, which the
note gets from $z\leq p^{h/2}$ and the definition of $s_0$, is obtained instead
from $p^8\leq K^2s^6\leq p^2s^6$ (\texttt{BowtieLeaf.le\_biased}).  The
monotonicity of $A_3$ on $[1/2,1]$, stated after Lemma~\ref{lem:scalar}, is not
formalized and is not needed.

\paragraph{Both symmetrisations are measure-preserving permutations.}
``Average over the automorphism group of the bowtie'' and ``by triangle
symmetry, the three choices of unweighted vertex give equal integrals'' are
discharged by \texttt{Methods/Peeling.integral\_assignment\_perm}, a general
lemma extracting the measure-preserving half of \texttt{homDensity\_iso} so
that it applies to an integrand which is \emph{not} a graph weight.  For the
outer orbit the three nontrivial automorphisms are given explicitly
(\texttt{aut2}, \texttt{aut3}, \texttt{aut4}, each checked by \texttt{decide});
for the triangle they are the two rotations.  Only transitivity on
$\{a,b,c,d\}$ is used, never the full group.  The two arithmetic--geometric
mean steps are Mathlib's
\texttt{Real.geom\_mean\_le\_arith\_mean\_weighted} at four equal weights $1/4$
and at three equal weights $1/3$; the second one carries the exponent
bookkeeping $\bigl(d^{1/4}d^{1/4}\bigr)^{1/3}\rightsquigarrow d^{1/6}$
(\texttt{BowtieLeaf.prod\_pair\_rootDegree}).

\paragraph{The factorisation \eqref{eq:square} is an identity, not an
estimate.}
\texttt{BowtieLeaf.integral\_outer\_factor} proves
$\int_{\Omega^5}=\int_\Omega R_\alpha^2$ exactly, by peeling the two triangles
into disjoint variable blocks; the Cauchy--Schwarz step
$\int R^2\geq(\int R)^2$ is separate
(\texttt{BowtieLeaf.sq\_integral\_pairTri\_le}, the project's
\texttt{BookTail.sq\_integral\_mul\_le} against the constant $1$).  The
accepted triangle input \eqref{eq:triangle-input} is
\texttt{K4Tail.goodman}, applied on $(\Omega,\nu)$ with no hypothesis on the
biased density.

\end{document}
